\section{Bitácora de actividades}\label{sec:bitacora}

Esta sección registra en orden cronológico cada actividad realizada: qué se hizo, cómo se hizo
y qué se obtuvo.

% Plantilla para nuevas entradas:
% \subsection{AAAA-MM-DD --- Título de la actividad}
% \textbf{Objetivo:} ...
% \textbf{Procedimiento:} ...
% \textbf{Resultado:} ...
% \textbf{Observaciones / problemas:} ...

\subsection{2026-09-26 --- Definición del problema}
\textbf{Objetivo:} delimitar la pregunta de investigación y el alcance de la fase~1.\\
\textbf{Procedimiento:} se definieron los periodos de análisis, las unidades temporal y territorial,
las variables objetivo candidatas y las hipótesis (Sección~\ref{sec:problema}).\\
\textbf{Resultado:} documentos \archivo{CONTEXTO.md} y \archivo{PENDIENTES.md} y este informe.

\subsection{2026-09-26 --- Identificación y verificación de fuentes}
\textbf{Objetivo:} localizar los conjuntos de datos que cubren las tres dimensiones y la pandemia.\\
\textbf{Procedimiento:} se consultó el catálogo del Portal de Datos Abiertos de la CDMX a través de
su API CKAN (\archivo{/api/3/action/package_search} y \archivo{package_show}). Se buscaron los
términos \emph{afluencia}, \emph{carpetas}, \emph{ecobici}, \emph{hechos de tránsito}, \emph{covid},
\emph{víctimas}, \emph{911}, \emph{metrobús}, \emph{empleo}, \emph{unidades económicas},
\emph{ocupación hotelera} y \emph{calidad del aire}. Para cada conjunto se revisó la lista de
archivos (recursos) que ofrece.\\
\textbf{Resultado:} se verificaron 20 conjuntos del portal (Sección~\ref{sec:fuentes}). Además:
\begin{itemize}
  \item Las carpetas de investigación de la FGJ se publican como CSV anuales de 2016 a 2024, más un
        archivo acumulado, en \url{https://archivo.datos.cdmx.gob.mx/FGJ/carpetas/}.
  \item Las llamadas al 911 se publican en archivos semestrales desde 2019-S1.
  \item El conjunto \enquote{Datos de movilidad (histórico COVID-19)} contiene índices relativos, no
        valores absolutos, por lo que no sirve para comparar con el periodo anterior a 2020.
\end{itemize}
\textbf{Observaciones:} el semáforo epidemiológico no está publicado como conjunto de datos y deberá
construirse a mano.

\subsection{2026-09-26 --- Inspección inicial de la afluencia del Metro}
\textbf{Objetivo:} comprobar el formato y la cobertura de la fuente principal de movilidad.\\
\textbf{Procedimiento:} se descargó el recurso \enquote{Afluencia Diaria del Metro (Simple)} y se
revisaron su encabezado, número de filas y rango de fechas con Python (módulo \archivo{csv}).\\
\textbf{Resultado:}
\begin{itemize}
  \item Columnas: \archivo{fecha, anio, mes, linea, estacion, afluencia}.
  \item 1,180,920 registros; cada registro corresponde a un día, una línea y una estación.
  \item Cobertura: del 2010-01-01 al 2026-07-31. La línea base anterior a la pandemia es suficiente.
\end{itemize}
\textbf{Observaciones / problemas:} los nombres de las estaciones presentan \emph{mojibake}
(\texttt{Isabel la Católica}). Se registró como tarea de limpieza.

\subsection{2026-09-26 --- Entorno de Python y notebook de contexto}
\textbf{Objetivo:} contar con un notebook que explique la problemática, lo que se analizará y los
conjuntos de datos, y que sirva como primera exploración reproducible.\\
\textbf{Procedimiento:}
\begin{itemize}
  \item Se creó el entorno virtual \archivo{.venv} con \archivo{pandas}, \archivo{requests},
        \archivo{matplotlib}, \archivo{ftfy}, \archivo{pyarrow} y \archivo{jupyter}.
  \item Se elaboró el notebook \archivo{notebooks/00_contexto_y_datasets.ipynb}. Contiene la
        problemática, la pregunta, las dimensiones y variables, los periodos, las hipótesis y un
        catálogo de 29 conjuntos de datos.
  \item El notebook verifica en vivo, mediante la API CKAN, que cada conjunto del portal existe, y
        muestra su fecha de actualización y su número de archivos.
  \item Se descargó la afluencia diaria del Metro en \archivo{data/raw/metro/} (58~MB) para hacer
        una primera exploración.
\end{itemize}
\textbf{Resultado:}
\begin{itemize}
  \item La afluencia mensual del Metro tocó fondo en \textbf{mayo de 2020}, con un índice de 27
        (2019 = 100), es decir, una caída del 73\,\%. En julio de 2026 el índice es de 78.1: la
        afluencia \textbf{aún no recupera el nivel previo a la pandemia} (Figura~\ref{fig:metro}).
  \item Promedio anual del índice: 2020 = 56.1, 2021 = 49.8, 2022 = 64.6, 2023 = 69.9,
        2024 = 73.5 y 2025 = 78.7.
  \item La caída aislada de septiembre de 2017 corresponde al sismo del 19-S, un ejemplo de choque
        puntual frente al choque prolongado de la pandemia.
\end{itemize}
\textbf{Problemas de calidad detectados:}
\begin{enumerate}
  \item \emph{Mojibake} en 57 de 168 nombres de estación. Se reparó con \archivo{ftfy}; tras la
        reparación quedan 165 estaciones únicas, porque había variantes duplicadas.
  \item Los \textbf{nombres de línea son inconsistentes}: entre 2021 y 2023 aparece
        \texttt{LÃnea 1} y en los demás años \archivo{Linea 1}. Por eso hay 24 valores en lugar
        de 12 líneas, y cada línea queda partida en dos series. Se normalizó a \archivo{Línea N}.
  \item Los cierres de la Línea~12 (2021--2023) y de la Línea~1 (2022--2024) reducen la afluencia
        por causas ajenas a la demanda. Deben marcarse en la limpieza.
\end{enumerate}

El análisis detallado de estos resultados está en la Sección~\ref{sec:exploracion}.


\subsection{2026-09-26 --- Ampliación del informe con los resultados del notebook}
\textbf{Objetivo:} que el informe contenga la misma información que el notebook de contexto.\\
\textbf{Procedimiento:} se añadieron varios elementos al informe:
\begin{itemize}
  \item La tabla de problemas de los datos y la tabla de dimensiones y variables (Sección~\ref{sec:problema}).
  \item La verificación en vivo de los 20 conjuntos del portal (Tabla~\ref{tab:verificacion}).
  \item Una sección nueva con la exploración de la afluencia del Metro (Sección~\ref{sec:exploracion}):
        perfil de las columnas, problemas de calidad, figura, índice $2019 = 100$ y afluencia por línea.
\end{itemize}
\textbf{Resultado:} sin contar los cierres de las Líneas~1 y 12, la afluencia de 2025 equivale al
83.2\,\% de la de 2019; con ellos, al 78.7\,\%. Ninguna línea ha recuperado su nivel previo a la
pandemia.

\subsection{2026-09-26 --- Selección de conjuntos de datos}
\textbf{Objetivo:} reducir el inventario de 29 fuentes a un conjunto manejable que responda la pregunta.\\
\textbf{Procedimiento:} se aplicaron cuatro criterios: línea base anterior a 2020, cobertura hasta
2024--2025, granularidad mensual o más fina, y descarga sin registro. Se revisaron los encabezados
de los CSV de los candidatos directamente en el portal.\\
\textbf{Resultado:} se seleccionaron 14 conjuntos (Tabla~\ref{tab:seleccion}); el resto quedó como opcional.
La revisión de los encabezados mostró que:
\begin{itemize}
  \item La ocupación hotelera empieza en octubre de 2018, por lo que su línea base es corta.
  \item El Metrobús tiene una serie simple desde 2005 y una desglosada por tipo de pago desde 2021.
  \item Ecobici ofrece dos formatos: diario por género y horario por rango de edad (desde 2010).
  \item La serie de casos COVID de la CDMX empieza el 5 de marzo de 2020.
\end{itemize}

\subsection{2026-09-26 --- Descarga de los conjuntos seleccionados}
\textbf{Objetivo:} obtener los 14 conjuntos seleccionados de forma reproducible y trazable.\\
\textbf{Procedimiento:} se escribieron seis scripts de descarga (Sección~\ref{sec:descarga}). Se
descargó primero el portal de la CDMX (unos 1.5~GB), luego INEGI y por último el IMSS, filtrado a
la CDMX. Antes de descargar se consultó el tamaño de cada archivo.\\
\textbf{Resultado:} se descargaron 13 de los 14 conjuntos; el SESNSP queda como descarga manual. Todo
está registrado en \archivo{data/raw/MANIFEST.csv}. La cobertura real aparece en la Tabla~\ref{tab:cobertura}.\\
\textbf{Observaciones / problemas:}
\begin{itemize}
  \item SharePoint (SESNSP) identifica el script como bot y devuelve una página de error. No se intentó
        evadir el bloqueo; el archivo se descargará a mano.
  \item El servidor del IMSS envía los archivos comprimidos y en codificaciones distintas según el año.
        El script se ajustó para filtrar sobre bytes.
  \item Los nombres de archivo perdían letras acentuadas (\archivo{metrob_s}). Se corrigió la
        normalización y se renombraron los archivos en el MANIFEST.
  \item Se encontró que la serie comparable de hechos de tránsito no cubre la pandemia. Se agregó la
        serie ampliada.
\end{itemize}

\subsection{2026-09-26 --- Limpieza de las carpetas FGJ y del Metrobús}
\textbf{Objetivo:} obtener las primeras dos fuentes limpias y listas para integrarse.\\
\textbf{Procedimiento:}
\begin{itemize}
  \item Se creó el catálogo de alcaldías a partir del Marco Geoestadístico y funciones comunes de
        normalización (\archivo{src/clean/catalogos.py}).
  \item Se exploraron las columnas, las categorías, las fechas y los duplicados de la FGJ, y se leyeron
        sus tres notas metodológicas.
  \item Se definieron 22 grupos de delito estables (\archivo{grupos_delito.py}).
  \item Se escribieron los scripts \archivo{fgj.py} y \archivo{metrobus.py} (Sección~\ref{sec:limpieza}).
\end{itemize}
\textbf{Resultado:}
\begin{itemize}
  \item FGJ: 1,937,812 carpetas limpias, el 97.5\,\% del original.
  \item Metrobús: 7 líneas unificadas y 2 errores anulados.
  \item Aparecen primeros indicios a favor de H1 y H2, y un contraste Metro--Metrobús no previsto.
\end{itemize}
\textbf{Observaciones / problemas:}
\begin{itemize}
  \item La FGJ no publica un identificador de carpeta, así que los duplicados no se pueden confirmar;
        se marcaron en vez de borrarse.
  \item El robo a transeúnte ya bajaba antes de 2020, así que el análisis tendrá que controlar la
        tendencia previa.
  \item El Metrobús publicó 57 valores con decimales en 2025, probablemente estimaciones.
\end{itemize}

\subsection{2026-09-26 --- Fin de la descarga del IMSS}
\textbf{Resultado:} se descargaron 128 meses, de enero de 2016 a agosto de 2026. Se conservaron solo
las filas de la CDMX (entre 430 y 550 mil por mes, de 3.6 a 4.7 millones nacionales). El total
ocupa 717~MB comprimidos, frente a unos 45~GB de los archivos nacionales completos.\\
\textbf{Observaciones / problemas:} la descarga se detuvo en enero de 2025 por una línea mal formada
(con menos columnas) en el archivo del IMSS. El script se ajustó para omitir y contar ese tipo de
líneas, y la descarga se reanudó desde ese mes. En total se omitió 1 línea.

\subsection{2026-09-26 --- Notebook de limpieza}
\textbf{Objetivo:} mostrar en un notebook los datos limpios y las decisiones de limpieza.\\
\textbf{Procedimiento:} se inició el notebook \archivo{notebooks/01_limpieza_datasets.ipynb}. Primero
se trabajó el semáforo: se filtró la CDMX, se pasó a formato largo y se guardó en
\archivo{data/interim/semaforo_cdmx_semana.csv}. Después se agregaron secciones con los resultados de
los scripts de limpieza de la FGJ y del Metrobús. Cada sección incluye el reporte de pasos, la carga de
los CSV limpios, la revisión de faltantes y gráficas.\\
\textbf{Observaciones:} al revisar el semáforo en formato largo se encontró que, además de la semana
2022-W01, falta la 2020-W53 (2020 tuvo 53 semanas ISO). La revisión con \archivo{isna()} daba cero
vacíos porque las semanas faltantes no aparecen como filas. Se detectaron comparando contra el
calendario completo de semanas.

\subsection{2026-09-26 --- Publicación del proyecto en GitHub}
\textbf{Problema:} el primer commit incluía todos los datos crudos (3~GB) y seis archivos de más de
100~MB, así que GitHub rechazaba el envío.\\
\textbf{Solución:} se agregó un \archivo{.gitignore} que excluye \archivo{data/raw/} y la base de
carpetas de la FGJ a nivel de registro (495~MB). Se versionan el código, los notebooks, el informe, el
\archivo{MANIFEST.csv} y los datos limpios agregados. El commit quedó en 6~MB. Los datos excluidos se
regeneran con los scripts de descarga y limpieza, y el MANIFEST permite verificar con SHA-256 que son
idénticos a los originales. Las instrucciones están en el \archivo{README.md}.

\subsection{2026-09-26 --- Limpieza del IMSS y criterios de imputación}
\textbf{Procedimiento:} se revisó el diccionario oficial del IMSS. Se definieron las columnas que se
conservan y la estructura de la tabla (mes por subdelegación). El código se escribió por pasos en el
notebook 01 y se ejecutó ahí (Sección~\ref{sec:limpieza}).\\
\textbf{Criterio de imputación:} se decidió imputar solo cuando el dato existió pero no se registró y
el hueco es corto. Casos con regresión: el desempleo de la ENOE en 2020-T2 y el backcasting de la
ocupación hotelera 2016--2018. Casos con métodos simples: el semáforo, dos días del Metrobús y los
meses finales de la FGJ. No se imputan los huecos largos (911 y series que terminan en 2024) ni los
datos que no existieron (cierres de líneas, personas sin empleo).\\
\textbf{Observaciones:} el rango de UMA pierde utilidad desde 2023 por el aumento del salario mínimo;
se usará el salario promedio deflactado.

\subsection{2026-09-26 --- Selección final de ocho conjuntos}
\textbf{Decisión:} se redujo la selección de 14 a 8 conjuntos (Sección~\ref{sec:seleccion_final}). Se
eligieron los hechos de tránsito en lugar de Ecobici porque complementan las debilidades del Metro:
tránsito vehicular y datos por alcaldía. Se descartaron el 911, la ENOE, el SESNSP y el DENUE por su
costo de limpieza, su cobertura o su redundancia.\\
\textbf{Revisión de la ocupación hotelera:} no tiene vacíos, meses faltantes ni valores fuera de rango
entre octubre de 2018 y julio de 2024. Los valores de abril a junio de 2020 (1.6--3.4\,\%) son reales:
corresponden al confinamiento. La precisión de los datos cambia desde febrero de 2020 (de 2 decimales
a muchos), lo que sugiere un cambio en el método de cálculo de la fuente.

\subsection{2026-09-26 --- Limpieza de la ocupación hotelera}
\textbf{Procedimiento:} en el notebook 01 se revisaron los tres tipos de faltantes; no se encontró
ninguno. Se redondearon los valores a dos decimales, se pasó la fecha al primer día del mes y se
renombró la columna (Sección~\ref{sec:limpieza}).\\
\textbf{Resultado:} \archivo{hoteles_mes.csv}, con 70 meses. Promedio de 67.7\,\% en 2019 y de
63.6\,\% en 2023.

\subsection{2026-09-26 --- Limpieza de los hechos de tránsito}
\textbf{Procedimiento:} se revisaron los folios, los duplicados y los nombres de alcaldía antes de
agregar. En el notebook 01 se unificaron las fechas, se unieron las dos series, se asignó la clave
INEGI y se agregaron los datos por mes, alcaldía y tipo de accidente. Cada error encontrado se
registró en una tabla dentro de la sección.\\
\textbf{Resultado:} \archivo{transito_alcaldia_mes.csv}, con 8,064 filas y 164,706 accidentes.\\
\textbf{Observaciones:} se agregó la errata \enquote{CUAHUTEMOC} al catálogo común de alcaldías
(\archivo{src/clean/catalogos.py}).

\subsection{2026-09-26 --- Visualización de hechos de tránsito y hoteles}
\textbf{Procedimiento:} se agregaron al notebook 01 cuatro gráficas de los hechos de tránsito (tipos
por mes, peatones y ciclistas, mapa de calor por alcaldía y fallecidos por año) y se guardó la de
ocupación hotelera. Las cinco se incorporaron al informe.\\
\textbf{Problemas detectados al graficar:} un corte de registro en las caídas de ciclista en febrero
de 2020, previo al confinamiento, y un registro incompleto en el primer semestre de 2018. Ambos se
agregaron a la tabla de errores del notebook. Se corrigió una interpretación preliminar que atribuía
el aumento de las caídas de ciclista a un mayor uso de la bicicleta.

\subsection{2026-09-27 --- Limpieza de la afluencia del Metro}
\textbf{Objetivo:} dejar lista la serie principal de movilidad y separar el efecto de la pandemia del
de los cierres de líneas.\\
\textbf{Procedimiento:} primero se exploró la afluencia por año con los datos crudos. Después, en el
notebook 01 se repararon la codificación y los nombres de línea y de estación, se corrigió el duplicado
de Oceanía, se validó el calendario y se clasificaron los ceros en cuatro estados. Luego se revisaron
los atípicos y se marcaron los cierres de las Líneas~1 y 12. Cada error quedó en una tabla al final de
la sección (Tabla~\ref{tab:limpieza_metro}).\\
\textbf{Resultado:} \archivo{metro_linea_dia}, \archivo{metro_estacion_mes} y \archivo{metro_mes} en
\archivo{data/interim/}, y tres gráficas.\\
\textbf{Observaciones:} la suma anual con los datos crudos no permite comparar años, porque 2017
pierde 8 días sin conteo (19-S) y 2026 solo llega a julio. Se usa el promedio diario. Sin las
Líneas~1 y 12, el Metro solo recuperó el 83\,\% de 2019.

\subsection{2026-09-27 --- Tabla diaria de hechos de tránsito}
\textbf{Objetivo:} tener los hechos de tránsito también por día, por si el análisis lo requiere.\\
\textbf{Procedimiento:} en el notebook 01 se agregaron los accidentes limpios por día, alcaldía y
tipo sobre la cuadrícula completa, con 0 donde no hubo accidentes. Igual que con el Metro, se revisó
el total diario contra su mediana móvil.\\
\textbf{Resultado:} \archivo{transito_alcaldia_dia} (245,472 filas), que cuadra mes por mes con la
tabla mensual.\\
\textbf{Observaciones:} no falta ningún día, pero siete días tienen un registro anormalmente bajo;
se marcaron con \archivo{registro\_bajo} y se agregaron a la tabla de errores.

\subsection{2026-09-27 --- Nombre de alcaldía en la FGJ y limpieza de casos COVID}
\textbf{Procedimiento:} se modificó \archivo{src/clean/fgj.py} para que la tabla
\archivo{fgj_alcaldia_mes} incluya el nombre de la alcaldía además de la clave, y se volvió a
ejecutar. Se verificó que el resto de la tabla no cambió. Después se limpiaron los casos positivos de
COVID-19 en el notebook 01: se completó el calendario, se revisó la consistencia entre columnas, se
analizó el efecto del día de la semana y se generaron las tablas diaria, semanal y mensual.\\
\textbf{Resultado:} \archivo{covid_dia}, \archivo{covid_semana} y \archivo{covid_mes}, con
1,617,973 positivos de residentes, y la gráfica semanal.\\
\textbf{Observaciones:} solo se encontró una fila inconsistente. La semana de los casos usa la misma
clave ISO que el semáforo, así que ambas fuentes se pueden unir directamente.

\subsection{2026-10-02 --- Corrección de formato del informe y limpieza explícita en el notebook}
\textbf{Problema:} en el PDF los separadores de miles se veían como espacios (por ejemplo,
\enquote{1 180 920}) en lugar de comas, y en el notebook 01 no se veía cómo se limpiaron todas las
fuentes: la FGJ y el Metrobús solo mostraban el reporte de sus scripts, y el semáforo y el IMSS tenían
código sin explicación.\\
\textbf{Procedimiento:} se cambiaron todos los separadores de miles del informe a coma (1,180,920) y se
corrigió la escritura \archivo{Línea N}. En el notebook 01 se agregó un índice de secciones y se
explicó cada paso del semáforo (incluida la detección de las semanas faltantes contra el calendario
completo) y del IMSS. Para la FGJ y el Metrobús se agregaron celdas que repiten a la vista cada paso de
sus scripts, con el número de registros afectados, y al final comprueban que el resultado es idéntico
a lo que guardó el script. También se corrigió la numeración de los pasos de la ocupación hotelera.\\
\textbf{Resultado:} el notebook se ejecutó completo sin errores y las comprobaciones contra los datos
del script dieron resultados idénticos.

\subsection{2026-10-02 --- Verificación de cifras, semáforo completo y versión en Word}
\textbf{Procedimiento:} se recalcularon desde los datos limpios todas las cifras que cita el informe
(Metro, Metrobús, FGJ, IMSS, ocupación hotelera, hechos de tránsito y casos COVID).\\
\textbf{Resultado:} casi todas coincidieron. Se corrigieron cinco puntos:
\begin{itemize}
  \item FGJ: el percentil 90 del retraso de la denuncia es de 56 días, no de 70. El valor anterior
        incluía hechos anteriores a 2016, que después se descartan.
  \item FGJ: los duplicados no tienen la hora de inicio vacía. Todos tienen \archivo{00:00:00}, un
        valor de relleno.
  \item IMSS: diciembre de 2020 no es el mínimo. El punto más bajo fue marzo de 2021 (3,213,626
        puestos). Se completaron en la tabla las mujeres y el salario de ambos meses.
  \item Hechos de tránsito: se agregaron las excepciones al patrón por alcaldía (Milpa Alta y
        Venustiano Carranza).
  \item Metro y COVID: se agregó la Línea~B a la lista de líneas rezagadas, y se aclaró que los 323
        días de la tabla de COVID son fines de semana.
\end{itemize}
Además, se rellenaron las semanas 2020-W53 y 2022-W01 del semáforo con el color de sus semanas
vecinas, que era el mismo antes y después, y se marcaron con \archivo{imputado}. Por último, el informe
se exportó a Word (\archivo{informe/informe.docx}) con el script \archivo{informe/a_word.py}, que usa
pandoc y se puede volver a correr después de cualquier cambio.
